\chapter*{Anexo A}
\addcontentsline{toc}{chapter}{Anexo A: Fuerzas de diseño por estación}

\setcounter{table}{0}
\renewcommand{\thetable}{A.\arabic{table}}
\setcounter{figure}{0}
\renewcommand{\thefigure}{A.\arabic{figure}}

\section*{Fuerzas de diseño por estación para ingreso al modelo FEM}
\label{anexo:fuerzas}

Las tablas de este anexo entregan, para cada una de las $21$ estaciones de la pala, la
fuerza aerodinámica que la simulación registra y la fuerza puntual con que esa estación
se carga en el modelo de elementos finitos. Se presentan cuatro tablas, una por cada
combinación de caso de carga último y sentido de la componente normal, según se explica
en la Sección~\ref{sec:comparacion_dlc}.

Las columnas se interpretan como sigue.

\begin{itemize}
    \item \textbf{Est.} numera las estaciones desde la base de la pala hacia su extremo
          superior, y \textbf{z} es la cota del centroide de cada una sobre la base, que
          es el punto en que la carga remota debe aplicarse en el modelo.
    \item \textbf{$\Delta s$} es la longitud de pala que representa la estación. Los
          $21$ valores suman por construcción los $23{,}000$~m de envergadura.
    \item \textbf{$F_{sim,n}$} y \textbf{$F_{sim,t}$} son la fuerza aerodinámica por
          unidad de longitud que entrega la simulación, en N/m.
    \item \textbf{$F_{d,n}$} y \textbf{$F_{d,t}$} son la fuerza puntual de diseño, en N,
          que resulta de multiplicar la anterior por $\Delta s$ y por el factor de
          escalamiento del caso.
\end{itemize}

Los valores conservan su signo, y el signo es el que fija el sistema local de la pala de
la Figura~\ref{fig:ejes_locales}, determinado por medición en la
Sección~\ref{sec:ejes_locales_res}. La componente normal $F_n$ es \textbf{positiva hacia
el eje del rotor} y negativa hacia afuera, esto es, en el sentido de la fuerza
centrífuga. La componente tangencial $F_t$ es \textbf{positiva en el sentido del
movimiento de la pala}, que en un H-rotor apunta hacia el borde de ataque. Respetar
ambos signos al ingresar las cargas es indispensable, ya que de ellos depende si la
solicitación aerodinámica se suma o se resta a la inercial.

\begin{table}[H]
    \centering
    \caption{Fuerzas por estación para el DLC 1.1, carga normal hacia el eje del rotor.}
    \label{tab:anexo_fuerzas_dlc11}
    \begin{tabular}{ccrrrrr}
        \hline
        \textbf{Est.} & $\mathbf{z}$ [m] & $\mathbf{\Delta s}$ [m] & $\mathbf{F_{sim,n}}$ [N/m] & $\mathbf{F_{sim,t}}$ [N/m] & $\mathbf{F_{d,n}}$ [N] & $\mathbf{F_{d,t}}$ [N] \\
        \hline
        1 & 0{,}064 & 0{,}1284 & 821 & 34 & 158 & 6 \\
        2 & 0{,}320 & 0{,}3825 & 2556 & 380 & 1461 & 217 \\
        3 & 0{,}825 & 0{,}6279 & 3217 & 587 & 3019 & 551 \\
        4 & 1{,}569 & 0{,}8594 & 3604 & 727 & 4628 & 934 \\
        5 & 2{,}534 & 1{,}0716 & 3691 & 732 & 5911 & 1172 \\
        6 & 3{,}700 & 1{,}2600 & 3783 & 740 & 7124 & 1393 \\
        7 & 5{,}040 & 1{,}4201 & 3863 & 757 & 8198 & 1606 \\
        8 & 6{,}524 & 1{,}5486 & 3853 & 760 & 8916 & 1760 \\
        9 & 8{,}120 & 1{,}6424 & 3825 & 762 & 9388 & 1871 \\
        10 & 9{,}791 & 1{,}6996 & 3911 & 798 & 9935 & 2027 \\
        11 & 11{,}500 & 1{,}7188 & 4205 & 942 & 10801 & 2421 \\
        12 & 13{,}209 & 1{,}6996 & 4520 & 1050 & 11480 & 2667 \\
        13 & 14{,}880 & 1{,}6424 & 5118 & 1361 & 12563 & 3340 \\
        14 & 16{,}476 & 1{,}5486 & 5641 & 1668 & 13055 & 3860 \\
        15 & 17{,}960 & 1{,}4201 & 5940 & 1448 & 12605 & 3072 \\
        16 & 19{,}300 & 1{,}2600 & 6021 & 1754 & 11336 & 3303 \\
        17 & 20{,}466 & 1{,}0716 & 6059 & 1760 & 9703 & 2819 \\
        18 & 21{,}431 & 0{,}8594 & 5677 & 1544 & 7291 & 1983 \\
        19 & 22{,}175 & 0{,}6279 & 5223 & 1292 & 4902 & 1212 \\
        20 & 22{,}680 & 0{,}3825 & 4615 & 1008 & 2638 & 576 \\
        21 & 22{,}936 & 0{,}1284 & 2032 & 201 & 390 & 39 \\
        \hline
    \end{tabular}
\end{table}

\begin{table}[H]
    \centering
    \caption{Fuerzas por estación para el DLC 1.1, carga normal hacia afuera.}
    \label{tab:anexo_fuerzas_dlc11_ext}
    \begin{tabular}{ccrrrrr}
        \hline
        \textbf{Est.} & $\mathbf{z}$ [m] & $\mathbf{\Delta s}$ [m] & $\mathbf{F_{sim,n}}$ [N/m] & $\mathbf{F_{sim,t}}$ [N/m] & $\mathbf{F_{d,n}}$ [N] & $\mathbf{F_{d,t}}$ [N] \\
        \hline
        1 & 0{,}064 & 0{,}1284 & $-$328 & $-$26 & $-$60 & $-$5 \\
        2 & 0{,}320 & 0{,}3825 & $-$1591 & 58 & $-$873 & 32 \\
        3 & 0{,}825 & 0{,}6279 & $-$2008 & 119 & $-$1809 & 107 \\
        4 & 1{,}569 & 0{,}8594 & $-$2183 & 147 & $-$2693 & 182 \\
        5 & 2{,}534 & 1{,}0716 & $-$2112 & 118 & $-$3249 & 182 \\
        6 & 3{,}700 & 1{,}2600 & $-$1919 & 77 & $-$3471 & 140 \\
        7 & 5{,}040 & 1{,}4201 & $-$1871 & 64 & $-$3812 & 131 \\
        8 & 6{,}524 & 1{,}5486 & $-$2134 & 101 & $-$4742 & 224 \\
        9 & 8{,}120 & 1{,}6424 & $-$2216 & 126 & $-$5224 & 296 \\
        10 & 9{,}791 & 1{,}6996 & $-$2595 & 196 & $-$6331 & 477 \\
        11 & 11{,}500 & 1{,}7188 & $-$3199 & 317 & $-$7891 & 782 \\
        12 & 13{,}209 & 1{,}6996 & $-$3373 & 335 & $-$8228 & 817 \\
        13 & 14{,}880 & 1{,}6424 & $-$3709 & 418 & $-$8741 & 986 \\
        14 & 16{,}476 & 1{,}5486 & $-$4239 & 573 & $-$9421 & 1273 \\
        15 & 17{,}960 & 1{,}4201 & $-$4885 & 811 & $-$9957 & 1653 \\
        16 & 19{,}300 & 1{,}2600 & $-$5122 & 920 & $-$9262 & 1664 \\
        17 & 20{,}466 & 1{,}0716 & $-$5156 & 928 & $-$7929 & 1427 \\
        18 & 21{,}431 & 0{,}8594 & $-$4762 & 752 & $-$5873 & 927 \\
        19 & 22{,}175 & 0{,}6279 & $-$4148 & 530 & $-$3738 & 477 \\
        20 & 22{,}680 & 0{,}3825 & $-$3407 & 328 & $-$1870 & 180 \\
        21 & 22{,}936 & 0{,}1284 & $-$1430 & 12 & $-$264 & 2 \\
        \hline
    \end{tabular}
\end{table}

\begin{table}[H]
    \centering
    \caption{Fuerzas por estación para el DLC 1.3, carga normal hacia el eje del rotor.}
    \label{tab:anexo_fuerzas_dlc13}
    \begin{tabular}{ccrrrrr}
        \hline
        \textbf{Est.} & $\mathbf{z}$ [m] & $\mathbf{\Delta s}$ [m] & $\mathbf{F_{sim,n}}$ [N/m] & $\mathbf{F_{sim,t}}$ [N/m] & $\mathbf{F_{d,n}}$ [N] & $\mathbf{F_{d,t}}$ [N] \\
        \hline
        1 & 0{,}064 & 0{,}1284 & 890 & 41 & 154 & 7 \\
        2 & 0{,}320 & 0{,}3825 & 2739 & 441 & 1414 & 228 \\
        3 & 0{,}825 & 0{,}6279 & 3280 & 614 & 2781 & 521 \\
        4 & 1{,}569 & 0{,}8594 & 3545 & 706 & 4113 & 819 \\
        5 & 2{,}534 & 1{,}0716 & 3858 & 817 & 5582 & 1182 \\
        6 & 3{,}700 & 1{,}2600 & 3995 & 859 & 6796 & 1460 \\
        7 & 5{,}040 & 1{,}4201 & 4108 & 929 & 7876 & 1780 \\
        8 & 6{,}524 & 1{,}5486 & 4114 & 888 & 8600 & 1855 \\
        9 & 8{,}120 & 1{,}6424 & 3956 & 862 & 8773 & 1910 \\
        10 & 9{,}791 & 1{,}6996 & 4180 & 958 & 9592 & 2199 \\
        11 & 11{,}500 & 1{,}7188 & 4115 & 947 & 9547 & 2196 \\
        12 & 13{,}209 & 1{,}6996 & 4601 & 1133 & 10556 & 2600 \\
        13 & 14{,}880 & 1{,}6424 & 4727 & 1195 & 10481 & 2650 \\
        14 & 16{,}476 & 1{,}5486 & 4924 & 1229 & 10294 & 2570 \\
        15 & 17{,}960 & 1{,}4201 & 5196 & 1346 & 9962 & 2581 \\
        16 & 19{,}300 & 1{,}2600 & 5501 & 1523 & 9356 & 2591 \\
        17 & 20{,}466 & 1{,}0716 & 5750 & 1545 & 8319 & 2236 \\
        18 & 21{,}431 & 0{,}8594 & 5757 & 1674 & 6679 & 1943 \\
        19 & 22{,}175 & 0{,}6279 & 5485 & 1508 & 4650 & 1279 \\
        20 & 22{,}680 & 0{,}3825 & 4809 & 1141 & 2483 & 589 \\
        21 & 22{,}936 & 0{,}1284 & 2213 & 248 & 384 & 43 \\
        \hline
    \end{tabular}
\end{table}

\begin{table}[H]
    \centering
    \caption{Fuerzas por estación para el DLC 1.3, carga normal hacia afuera.}
    \label{tab:anexo_fuerzas_dlc13_ext}
    \begin{tabular}{ccrrrrr}
        \hline
        \textbf{Est.} & $\mathbf{z}$ [m] & $\mathbf{\Delta s}$ [m] & $\mathbf{F_{sim,n}}$ [N/m] & $\mathbf{F_{sim,t}}$ [N/m] & $\mathbf{F_{d,n}}$ [N] & $\mathbf{F_{d,t}}$ [N] \\
        \hline
        1 & 0{,}064 & 0{,}1284 & $-$428 & $-$25 & $-$74 & $-$4 \\
        2 & 0{,}320 & 0{,}3825 & $-$1939 & 110 & $-$1001 & 57 \\
        3 & 0{,}825 & 0{,}6279 & $-$2466 & 207 & $-$2091 & 175 \\
        4 & 1{,}569 & 0{,}8594 & $-$2778 & 277 & $-$3223 & 322 \\
        5 & 2{,}534 & 1{,}0716 & $-$3111 & 359 & $-$4501 & 520 \\
        6 & 3{,}700 & 1{,}2600 & $-$3246 & 384 & $-$5522 & 653 \\
        7 & 5{,}040 & 1{,}4201 & $-$3245 & 375 & $-$6222 & 719 \\
        8 & 6{,}524 & 1{,}5486 & $-$3306 & 393 & $-$6912 & 821 \\
        9 & 8{,}120 & 1{,}6424 & $-$3448 & 450 & $-$7646 & 998 \\
        10 & 9{,}791 & 1{,}6996 & $-$3391 & 431 & $-$7781 & 990 \\
        11 & 11{,}500 & 1{,}7188 & $-$3566 & 495 & $-$8275 & 1148 \\
        12 & 13{,}209 & 1{,}6996 & $-$3870 & 561 & $-$8879 & 1287 \\
        13 & 14{,}880 & 1{,}6424 & $-$3931 & 564 & $-$8715 & 1249 \\
        14 & 16{,}476 & 1{,}5486 & $-$4049 & 584 & $-$8464 & 1222 \\
        15 & 17{,}960 & 1{,}4201 & $-$4293 & 678 & $-$8230 & 1301 \\
        16 & 19{,}300 & 1{,}2600 & $-$4485 & 775 & $-$7629 & 1319 \\
        17 & 20{,}466 & 1{,}0716 & $-$4685 & 851 & $-$6778 & 1232 \\
        18 & 21{,}431 & 0{,}8594 & $-$4742 & 838 & $-$5501 & 973 \\
        19 & 22{,}175 & 0{,}6279 & $-$4346 & 661 & $-$3684 & 561 \\
        20 & 22{,}680 & 0{,}3825 & $-$3594 & 402 & $-$1856 & 207 \\
        21 & 22{,}936 & 0{,}1284 & $-$1552 & 27 & $-$269 & 5 \\
        \hline
    \end{tabular}
\end{table}

